\begin{frame}{Tratamento no GNU Octave}
\begin{enumerate}
\item Recebemos o bloco, verificamos o status e salvamos as amostras em \texttt{.mat}.
\item Estimamos período e fase para alinhar a triangular simulada ao sinal capturado.
\item Comparamos a referência contínua, a versão quantizada e os códigos do ADC.
\item Calculamos os erros no tempo, o histograma de códigos e os espectros.
\end{enumerate}
\begin{block}{Comparação utilizada}
O erro bruto compara os códigos com a triangular ideal de 0 a 255. O erro residual remove ganho e offset ajustados. Ele ainda inclui efeitos do gerador, do ADC e do modelo.
\end{block}
\pblfonte{adc\_analyzer.m e módulos de src/; relatório técnico, seção 4.}
\end{frame}

\begin{frame}{Comparação no Tempo: Ensaio de 10 Hz}
\centering
\begin{tikzpicture}
\begin{axis}[pblplot,height=5cm,xlabel={Índice da amostra após trigger},ylabel={Código ADC},xmin=0,xmax=1750,ymin=-5,ymax=260,legend pos=north east]
\addplot[green!50!black,dashed,thick] table[x index=0,y index=2]{resultados/tempo_10hz.dat};
\addplot[blue,const plot] table[x index=0,y index=1]{resultados/tempo_10hz.dat};
\legend{Triangular ajustada,ADC}
\end{axis}
\end{tikzpicture}
\par\small $f=10{,}072$ Hz, $f_s=8196{,}72$ Hz e RMSE bruto de $0{,}393$ LSB.
\pblfonte{sample\_20261106\_102822\_177.mat; ajuste pelo adc\_analyzer.}
\end{frame}

\begin{frame}{Comparação no Domínio da Frequência}
\centering
\begin{tikzpicture}
\begin{axis}[pblplot,height=4.7cm,xlabel={Frequência (Hz)},ylabel={Magnitude (dBFS)},xmin=0,xmax=170,ymin=-105,ymax=5,legend pos=north east]
\addplot[green!50!black,dashed] table[x index=0,y index=3]{resultados/fft_10hz.dat};
\addplot[orange] table[x index=0,y index=2]{resultados/fft_10hz.dat};
\addplot[blue] table[x index=0,y index=1]{resultados/fft_10hz.dat};
\legend{Triangular contínua,Ideal quantizada,ADC}
\end{axis}
\end{tikzpicture}
\par\small As harmônicas ímpares fazem parte da triangular. A comparação usa remoção da média e janela de Hann para reduzir o vazamento espectral.
\pblfonte{Captura de 06/11, 10:28:22; computeSpectralMetrics.m.}
\end{frame}
